\phantomsection
\chapter*{Simgeler ve Kısaltmalar}
\markboth{Simgeler ve Kısaltmalar}{Simgeler ve Kısaltmalar}
\addcontentsline{toc}{chapter}{Simgeler ve Kısaltmalar}

Bu bölümde kitap boyunca sık kullanılan temel simgeler ve kısaltmalar toplu
olarak verilmiştir. Büyük harfli $\bar X$, $S^2$ veya $\hat\Theta$ gösterimi
çoğunlukla örneklem seçilmeden önce rassal olan tahmin ediciyi; küçük harfli
$\bar x$, $s^2$ veya $\hat\theta$ ise gözlenen örneklemden hesaplanan sayısal
tahmini belirtir. Uygulamalı bölümlerde bağlam açık olduğunda bu ayrım daha
yalın gösterimlerle ifade edilebilir.

\small
\begingroup
\renewcommand{\arraystretch}{1.12}

\Needspace{18\baselineskip}
\noindent{\sffamily\bfseries Temel büyüklükler ve olasılık simgeleri}\par\smallskip
\begin{center}
\begin{tabular}{@{}>{\RaggedRight\arraybackslash}p{0.23\textwidth}>{\RaggedRight\arraybackslash}p{0.69\textwidth}@{}}
\toprule
\textbf{Simge} & \textbf{Anlamı} \\
\midrule
$n$ & Örneklem hacmi veya geçerli gözlem sayısı \\
$N$ & Evren hacmi ya da bağlama göre toplam gözlem sayısı \\
$x_i,\ y_i$ & $i$'inci birimin gözlenen değerleri \\
$X,Y$ & Rassal değişkenler \\
$P(A)$ & $A$ olayının olasılığı \\
$P(A\mid B)$ & $B$ gerçekleşmişken $A$ olayının koşullu olasılığı \\
$E[X]$ & $X$ rassal değişkeninin beklenen değeri \\
$\operatorname{Var}(X)$ veya $V(X)$ & $X$ rassal değişkeninin varyansı \\
$\operatorname{Cov}(X,Y)$ & $X$ ve $Y$ arasındaki kovaryans \\
$\mu$ & Evren ortalaması \\
$\sigma^2,\ \sigma$ & Evren varyansı ve evren standart sapması \\
\bottomrule
\end{tabular}
\end{center}

\Needspace{18\baselineskip}
\noindent{\sffamily\bfseries Örneklem ve tahmin simgeleri}\par\smallskip
\begin{center}
\begin{tabular}{@{}>{\RaggedRight\arraybackslash}p{0.23\textwidth}>{\RaggedRight\arraybackslash}p{0.69\textwidth}@{}}
\toprule
\textbf{Simge} & \textbf{Anlamı} \\
\midrule
$\bar X,\ \bar x$ & Örneklem ortalaması tahmin edicisi ve gözlenen örneklem ortalaması \\
$S^2,\ s^2$ & Örneklem varyansı tahmin edicisi ve gözlenen örneklem varyansı \\
$s$ & Gözlenen örneklem standart sapması \\
$p,\ \hat p$ & Evren oranı ve örneklem oranı \\
$\theta,\ \hat\theta$ & Hedef parametre ve bu parametre için tahmin \\
$SE(\hat\theta)$ & Bir tahmin edicinin standart hatası \\
$Q_1,\ Q_3$ & Birinci ve üçüncü çeyrek \\
$IQR=Q_3-Q_1$ & Çeyrekler arası açıklık \\
$Md$ & Medyan \\
$f_j,\ F_j$ & $j$'inci sınıfın frekansı ve kümülatif frekansı \\
$m_j$ & $j$'inci sınıf aralığının orta noktası \\
$\pi_i$ & $i$'inci birimin örnekleme dahil olma olasılığı \\
\bottomrule
\end{tabular}
\end{center}

\Needspace{19\baselineskip}
\noindent{\sffamily\bfseries İstatistiksel çıkarım simgeleri}\par\smallskip
\begin{center}
\begin{tabular}{@{}>{\RaggedRight\arraybackslash}p{0.23\textwidth}>{\RaggedRight\arraybackslash}p{0.69\textwidth}@{}}
\toprule
\textbf{Simge} & \textbf{Anlamı} \\
\midrule
$H_0,\ H_1$ & Sıfır hipotezi ve alternatif hipotez \\
$\alpha$ & Anlamlılık düzeyi veya Tip I hata olasılığı \\
$\beta$ & Tip II hata olasılığı \\
$1-\beta$ & Test gücü \\
$p$-değeri & $H_0$ ve test varsayımları altında gözlenen kadar veya daha aşırı sonuç elde etme olasılığı \\
$sd$ & Serbestlik derecesi \\
$z,\ t,\ \chi^2,\ F$ & İlgili yöntemlerde kullanılan test istatistikleri \\
$z^*,\ t^*$ & Güven aralığı veya test için seçilen kritik değerler \\
$z_\alpha,\ t_{\alpha,\nu}$ & Sağ kuyruğunda $\alpha$ olasılık bırakan kritik değerler \\
$\chi^2_{\alpha,\nu}$ & $\nu$ serbestlik dereceli ki-kare dağılımının kritik değeri \\
$F_{\alpha,d_1,d_2}$ & $d_1$ ve $d_2$ serbestlik dereceli F dağılımının kritik değeri \\
\bottomrule
\end{tabular}
\end{center}

\Needspace{21\baselineskip}
\noindent{\sffamily\bfseries Korelasyon, regresyon ve ANOVA simgeleri}\par\smallskip
\begin{center}
\begin{tabular}{@{}>{\RaggedRight\arraybackslash}p{0.23\textwidth}>{\RaggedRight\arraybackslash}p{0.69\textwidth}@{}}
\toprule
\textbf{Simge} & \textbf{Anlamı} \\
\midrule
$r,\ \rho$ & Örneklem Pearson korelasyonu ve evren korelasyonu \\
$r_s$ & Spearman sıra korelasyonu \\
$R^2$ & Regresyon modelinin belirlilik katsayısı \\
$\beta_0,\ \beta_1$ & Evren regresyon doğrusunun kesme ve eğim katsayıları \\
$b_0,\ b_1$ & Örneklemden tahmin edilen kesme ve eğim katsayıları \\
$\widehat Y_i$ & $i$'inci birim için uydurulan veya tahmin edilen yanıt \\
$\varepsilon_i,\ e_i$ & Model hata terimi ve gözlenen artık \\
$SST,\ SSR,\ SSE$ & Regresyonda toplam, model ve hata kareler toplamı \\
$SS_A,\ SS_E,\ SS_T$ & ANOVA'da gruplar arası, hata ve toplam kareler toplamı \\
$MS_A,\ MS_E$ & ANOVA'da gruplar arası ve hata ortalama karesi \\
$MSE$ & Ortalama karesel hata; regresyon/ANOVA bağlamında hata ortalama karesi \\
$s_e$ & Regresyon modelinin artık standart hatası \\
\bottomrule
\end{tabular}
\end{center}

\Needspace{18\baselineskip}
\noindent{\sffamily\bfseries Etki büyüklüğü ve parametrik olmayan test simgeleri}\par\smallskip
\begin{center}
\begin{tabular}{@{}>{\RaggedRight\arraybackslash}p{0.23\textwidth}>{\RaggedRight\arraybackslash}p{0.69\textwidth}@{}}
\toprule
\textbf{Simge} & \textbf{Anlamı} \\
\midrule
$d,\ d_z$ & Cohen'in standartlaştırılmış ortalama farkı ve eşleştirilmiş farklara dayalı sürümü \\
$g$ & Küçük örneklem yanlılığı düzeltilmiş Hedges etki büyüklüğü \\
$\eta^2,\ \omega^2$ & ANOVA için eta-kare ve yanlılık düzeltmeli omega-kare \\
$V$ & Kategorik ilişki için Cramér $V$ etki büyüklüğü \\
$OR$ & Odds oranı \\
$U$ & Mann--Whitney test istatistiği \\
$W$ & Wilcoxon işaretli sıra test istatistiği veya sıra toplamı \\
$H$ & Kruskal--Wallis test istatistiği \\
$\varepsilon^2$ & Kruskal--Wallis için epsilon-kare etki büyüklüğü \\
$r_{rb}$ & Sıra-biserial korelasyon \\
\bottomrule
\end{tabular}
\end{center}

\Needspace{20\baselineskip}
\noindent{\sffamily\bfseries Temel kısaltmalar}\par\smallskip
\begin{center}
\begin{tabular}{@{}>{\RaggedRight\arraybackslash}p{0.23\textwidth}>{\RaggedRight\arraybackslash}p{0.69\textwidth}@{}}
\toprule
\textbf{Kısaltma} & \textbf{Açılımı} \\
\midrule
ANCOVA & Kovaryans analizi \\
ANOVA & Varyans analizi \\
GA & Güven aralığı \\
HSD & Tukey yöntemindeki \textit{honestly significant difference}; dürüst anlamlı fark \\
IQR & Çeyrekler arası açıklık \\
MLT & Merkezi Limit Teoremi \\
MSE & Ortalama karesel hata \\
PDR & Psikolojik Danışmanlık ve Rehberlik \\
SE & Standart hata için formüllerde kullanılan İngilizce kısaltma \\
SH & Standart hata \\
SPSS & \textit{Statistical Package for the Social Sciences}; istatistik yazılımı \\
SS & Standart sapma; ANOVA'daki alt indisli $SS$ gösterimleri kareler toplamını belirtir \\
\bottomrule
\end{tabular}
\end{center}

\endgroup
\normalsize